\Conclusion

В выпускной квалификационной работе разработано веб-приложение на основе RAG для работы с документацией PostgreSQL с учётом выбранной версии. Приложение ориентировано на инженерный сценарий: ответ должен быть связан с локально индексированным корпусом, выбранной основной версией PostgreSQL и списком источников.

В первой главе выполнен анализ предметной области, особенностей официальной документации PostgreSQL, проблемы учёта версии и существующих подходов поиска по технической документации. Сравнение показало, что обычный поиск, полнотекстовый поиск, поиск в локальном корпусе HTML-документации и вопросно-ответные системы без источников не дают одновременно связный ответ, контроль версии, локальную управляемость корпуса и трассируемость источников.

Во второй главе спроектировано веб-приложение: определены архитектура серверной части, маршруты API, сценарий обработки запроса, диаграмма процесса, модель данных, алгоритмы подготовки корпуса, разбиения документации на фрагменты, поиска, учёта версии, повторного ранжирования и проверки достаточности подтверждений.

В третьей главе описана программная реализация: FastAPI-приложение, PostgreSQL с pgvector, SQLAlchemy-модели, Alembic-миграции, Pydantic-схемы, подготовка официального и дополнительного корпусов, обработка пользовательского запроса, работа с Groq API, история запросов и пользовательский интерфейс. Приложение поддерживает версии PostgreSQL 16, 17 и 18, краткий, подробный и обучающий режимы ответа, локальную модель построения векторных представлений \texttt{BAAI/bge-m3} размерности 1024 и отображение источников.

В четвертой главе проведены тестирование и оценка соответствия требованиям. Автоматический запуск тестов практической части завершился результатом \texttt{91 passed, 3 skipped}. Пропущенные тесты относятся к интеграционным проверкам, для которых нужна отдельная тестовая база PostgreSQL+pgvector через \texttt{TEST\_DATABASE\_URL}. Проверка подтвердила работу маршрутов API, валидации версии и режима, контроля версии, повторного ранжирования, проверки достаточности подтверждений, пользовательского интерфейса и истории запросов.

Оценка результатов выполнена функционально и сценарно. Количественные метрики качества поиска и ранжирования в работе не рассчитывались, поэтому в выводах не заявляются Recall@K, MRR, nDCG, accuracy или абсолютная точность ответа. Ограничения разработанного решения связаны с полнотой локального корпуса, доступностью Groq API, заданным набором поддерживаемых версий и отсутствием отдельного эталонного набора для численной оценки качества ранжирования.

Поставленная цель достигнута: разработано и описано веб-приложение, которое принимает вопрос пользователя, версию PostgreSQL и режим ответа, извлекает релевантные фрагменты корпуса, учитывает выбранную версию, возвращает источники, сохраняет историю запросов и поддерживает безопасный отказ при недостатке подтверждённых данных.
